\section{实验}
\subsection{数据集与评估指标}
nuScenes 3D检测基准测试（Benchmark）\cite{nuscenes} 包含1000个多模态视频，每个视频时长约20秒，关键帧以2Hz频率标注。
每个样本包含来自6个相机的图像，覆盖360度全景视野。
视频分为700个训练、150个验证和150个测试。检测任务包含10个物体类别的140万个标注3D边界框。
nuScenes 使用地面平面上基于中心距离的四个不同阈值计算平均精度（Mean Average Precision, mAP），并包含五个真阳性指标：ATE（平移误差）、ASE（尺度误差）、AOE（朝向误差）、AVE（速度误差）和AAE（属性误差）。
此外，通过结合检测精度（mAP）与五个真阳性指标定义了 nuScenes 检测分数（NDS）。

\setlength{\tabcolsep}{3pt}
\setlength{\doublerulesep}{2\arrayrulewidth}
\renewcommand{\arraystretch}{1.1}
\begin{table*}[t]
    \caption{BEVFormer v2 与其他当前最优（SOTA）方法在 nuScenes 测试集上的3D检测结果。$^\dagger$ 表示 V2-99 \cite{Vovnet} 在深度估计（Depth Estimation）任务上使用额外数据 \cite{DD3D} 进行了预训练。$^\ddagger$ 表示使用 CBGS 的方法，该方法将1个epoch延长为4.5个epoch。本文仅训练 BEVFormer v2 24个epoch以与先前方法公平比较。}
    \label{table:sota_nuscenes}
    \centering

    \begin{tabular}{l|c|c|c|cc|ccccc}
        \toprule
        Method & Backbone & Epoch & Image Size & NDS & mAP & mATE & mASE & mAOE & mAVE & mAAE \\ 
        \midrule
        BEVFormer~\cite{bevformer} & V2-99$^\dagger$ & 24 & 900 $\times$ 1600 & 0.569 & 0.481 & 0.582 & 0.256 & 0.375 & 0.378 & 0.126 \\
        PolarFormer~\cite{polarformer} & V2-99$^\dagger$ & 24 & 900 $\times$ 1600 & 0.572 & 0.493 & 0.556 & 0.256 & 0.364 & 0.440 & 0.127 \\
        PETRv2~\cite{petrv2} & GLOM & 24 & 640 $\times$ 1600 & 0.582 & 0.490 & 0.561 & 0.243 & 0.361 & 0.343 & 0.120 \\
        BEVDepth~\cite{bevdepth} & V2-99$^\dagger$ & 90$^\ddagger$ & 640 $\times$ 1600 & 0.600 & 0.503 & 0.445 & 0.245 & 0.378 & 0.320 & 0.126 \\
        BEVStereo~\cite{bevstereo} & V2-99$^\dagger$ & 90$^\ddagger$ & 640 $\times$ 1600 & 0.610 & 0.525 & 0.431 & 0.246 & 0.358 & 0.357 & 0.138 \\
        \midrule
        \textbf{BEVFormer v2} & InternImage-B & 24 & 640 $\times$ 1600 & 0.620 & 0.540 & 0.488 & 0.251 & 0.335 & 0.302 & 0.122 \\
        \textbf{BEVFormer v2} & InternImage-XL & 24 & 640 $\times$ 1600 & 0.634 & 0.556 & 0.456 & 0.248 & 0.317 & 0.293 & 0.123 \\
        \bottomrule
    \end{tabular}

\end{table*}
\setlength{\tabcolsep}{3pt}
\setlength{\doublerulesep}{2\arrayrulewidth}
\renewcommand{\arraystretch}{1.1}
\begin{table*}[t]
    \caption{不同视图监督（View Supervision）组合的3D检测器在 nuScenes 验证集上的检测结果。所有模型均未使用时序（Temporal）信息训练。}
    \label{table:view_supervision}
    \centering

    \begin{tabular}{c|c|c|cc|cccccc}
        \toprule
        View Supervision & Backbone & Epoch & NDS & mAP & mATE & mASE & mAOE & mAVE & mAAE   \\ 
        \midrule 
        Perspective Only & ResNet-101 & 48 & 0.412 & 0.323 & 0.737 & 0.268 & 0.377 & 0.943 & 0.167 \\ 
        BEV Only & ResNet-101 & 48 & 0.426 & 0.355 & 0.751 & 0.275 & 0.429 & 0.847 & 0.215 \\ 
        Perspective \& BEV & ResNet-101 & 48 & 0.451 & 0.374 & 0.730 & 0.270 & 0.379 & 0.773 & 0.205 \\ 
        BEV \& BEV & ResNet-101 & 48 & 0.428 & 0.350 & 0.750 & 0.279 & 0.388 & 0.842 & 0.210 \\ 
        \midrule
        \bottomrule
    \end{tabular}
    
\end{table*}

\subsection{实验设置}
本文使用多种骨干网络进行实验：ResNet~\cite{ResNet}、DLA~\cite{DLA}、VoVNet~\cite{Vovnet} 和 InternImage \cite{InternImage}。
所有骨干网络均使用在 COCO 数据集 \cite{COCO} 的2D检测任务上预训练的检查点初始化。
除本文修改外，遵循 BEVFormer~\cite{bevformer} 的默认设置构建BEV检测头。
在表 \ref{table:sota_nuscenes} 和表 \ref{table:trick} 中，BEV头采用新的时序编码器利用时序信息。其他实验中，采用仅使用当前帧的单帧版本，类似 BEVFormer-S~\cite{bevformer}。
对于透视3D检测头，采用 DD3D~\cite{DD3D} 的实现，带有相机感知深度参数化（Camera-aware Depth Parameterization）。
透视损失和BEV损失的权重设为 $\lambda_{bev}=\lambda_{pers}=1$。
使用 AdamW~\cite{AdamW} 优化器，基础学习率设为 4e-4。

\subsection{基准测试结果}
本文将提出的 BEVFormer v2 与现有当前最优（SOTA）的BEV检测器进行比较，包括 BEVFormer~\cite{bevformer}、PolarFormer~\cite{polarformer}、PETRv2~\cite{petrv2}、BEVDepth~\cite{bevdepth} 和 BEVStereo~\cite{bevstereo}。
表 \ref{table:sota_nuscenes} 报告了 nuScenes 测试集上的3D目标检测结果。
BEVFormer、PolarFormer、BEVDepth 和 BEVStereo 使用的 V2-99 \cite{Vovnet} 骨干网络已在额外数据的深度估计任务上预训练，然后通过 DD3D \cite{DD3D} 在 nuScenes 数据集 \cite{nuscenes} 上微调。
相反，本文采用的 InternImage \cite{InternImage} 骨干网络使用来自 COCO~\cite{COCO} 检测任务的检查点初始化，没有任何3D预训练。
InternImage-B 的参数量与 V2-99 相似，但更能体现现代图像骨干网络设计的进步。
可以观察到，使用 InternImage-B 骨干网络的 BEVFormer v2 优于所有现有方法，表明在透视监督下，不再需要单目3D任务预训练的骨干网络。
使用 InternImage-XL 的 BEVFormer v2 以 63.4\% NDS 和 55.6\% mAP 超越了 nuScenes 相机3D目标检测排行榜上的所有方法，比第二名 BEVStereo 高出 2.4\% NDS 和 3.1\% mAP。
这一显著提升揭示了释放现代图像骨干网络在BEV识别中的巨大潜力所带来的收益。

\setlength{\tabcolsep}{3pt}
\setlength{\doublerulesep}{2\arrayrulewidth}
\renewcommand{\arraystretch}{1.1}
\begin{table*}[ht]
    \caption{透视监督（Perspective Supervision）在不同2D图像骨干网络（Image Backbone）上的结果，在 nuScenes 验证集上评测。"BEV Only" 和 "Perspective \& BEV" 与表 \ref{table:view_supervision} 相同。所有骨干网络均使用 COCO~\cite{COCO} 预训练权重初始化，所有模型均未使用时序（Temporal）信息训练。}
    \label{table:large_backbone}
    \centering

    \begin{tabular}{c|c|c|cc|cccccc}
        \toprule
        Backbone & Epoch & View Supervision & NDS & mAP & mATE & mASE & mAOE & mAVE & mAAE   \\ 
        \midrule
        ResNet-50 & 48 & BEV Only & 0.400 & 0.327 & 0.795 & 0.277 & 0.479 &0.871 & 0.210 \\ 
        ResNet-50 & 48 & Perspective \& BEV & 0.428 & 0.349 & 0.750 & 0.276 & 0.424 & 0.817 & 0.193 \\
        \midrule
        DLA-34 & 48 & BEV Only & 0.403 & 0.338 & 0.772 & 0.279 & 0.483 & 0.919 & 0.206 \\ 
        DLA-34 & 48 & Perspective \& BEV & 0.435 & 0.358 & 0.742 & 0.274 & 0.431 & 0.801 & 0.186 \\  
  
        \midrule
        ResNet-101 & 48 & BEV Only & 0.426 & 0.355 & 0.751 & 0.275 & 0.429 & 0.847 & 0.215  \\ 
        ResNet-101 & 48 & Perspective \& BEV & 0.451 & 0.374 & 0.730 & 0.270 & 0.379 & 0.773 & 0.205  \\   
        \midrule
        VoVNet-99 & 48 & BEV Only & 0.441 & 0.367  & 0.734 & 0.271 & 0.402 & 0.815  & 0.205 \\ 
        VoVNet-99 & 48 & Perspective \& BEV & 0.467 & 0.396  & 0.709  & 0.274 & 0.368 & 0.768 & 0.196 \\ 
        \midrule
        InternImage-B & 48 & BEV Only & 0.455 & 0.398  & 0.712 & 0.283 & 0.411 & 0.826  & 0.204 \\ 
        InternImage-B & 48 & Perspective \& BEV & 0.485 & 0.417  & 0.696  & 0.275 & 0.354 & 0.734 & 0.182 \\ 
        \bottomrule
    \end{tabular}
    
\end{table*}
\setlength{\tabcolsep}{3pt}
\setlength{\doublerulesep}{2\arrayrulewidth}
\renewcommand{\arraystretch}{1.1}
\begin{table*}[t]
    \caption{仅使用 BEV 监督（BEV Supervision）与同时使用透视和 BEV 监督（Perspective \& BEV Supervision）的模型在不同训练轮数下的对比。模型在 nuScenes 验证集上评测。所有模型均未使用时序（Temporal）信息训练。}
    \label{table:long_schedule}
    \centering

    \begin{tabular}{c|c|c|cc|cccccc}
        \toprule
        View Supervision & Backbone & Epoch  & NDS & mAP & mATE & mASE & mAOE & mAVE & mAAE   \\ 
        \midrule 
        \multirow{3}{*}{BEV Only} & \multirow{3}{*}{ResNet-50} & 24 & 0.379 & 0.322 & 0.803 & 0.280 & 0.549 & 0.954 & 0.240 \\ 
         &  & 48 & 0.400 & 0.327 & 0.795 & 0.277 & 0.479 & 0.871 & 0.210 \\ 
         &  & 72 & 0.410 & 0.335 & 0.771 & 0.280 & 0.458 & 0.848 & 0.216 \\ 
        \midrule
        \multirow{3}{*}{Perspective \& BEV} & \multirow{3}{*}{ResNet-50} & 24 &  0.414 & 0.351 & 0.732 & 0.271 & 0.505 & 0.899 & 0.204 \\ 
         &  & 48 & 0.428 & 0.349 & 0.750 & 0.276 & 0.424 & 0.817 & 0.193 \\ 
         &  & 72 & 0.428 & 0.351 & 0.741 & 0.279 & 0.419 & 0.835 & 0.196 \\ 
        \midrule
        \bottomrule
    \end{tabular}
\end{table*}
%\vspace{+5pt}

\subsection{消融实验与分析}
\subsubsection{透视监督的有效性}
为验证透视监督的有效性，本文在表 \ref{table:view_supervision} 中比较了不同视图监督组合的3D检测器，包括：
(1) 透视与BEV（Perspective \& BEV）——本文提出的 BEVFormer v2，整合了透视头和BEV头的两阶段检测器。
(2) 仅透视（Perspective Only）——本文模型中的单阶段透视检测器。
(3) 仅BEV（BEV Only）——本文模型中不使用混合物体查询的单阶段BEV检测器。
(4) BEV与BEV（BEV \& BEV）——具有两个BEV头的两阶段检测器，即将本文模型中的透视头替换为另一个利用BEV特征生成混合物体查询提议的BEV头。

与仅透视检测器相比，仅BEV检测器通过利用多视角图像实现了更好的NDS和mAP，但其mATE和mAOE更高，表明了BEV监督的潜在问题。
本文的透视与BEV检测器实现了最佳性能，分别超过仅BEV检测器2.5\% NDS和1.9\% mAP。
具体而言，透视与BEV检测器的mATE、mAOE和mAVE显著低于仅BEV检测器。
这一显著提升主要来自以下两个方面：
(1) 在常规视觉任务上预训练的骨干网络无法捕获3D场景中物体的某些属性，包括深度、朝向和速度，而受透视监督指导的骨干网络能够提取这些属性的信息。
(2) 与固定物体查询集相比，本文的混合物体查询包含第一阶段预测作为参考点，帮助BEV头定位目标物体。
为进一步确保提升并非来自两阶段流程本身，本文引入BEV与BEV检测器进行比较。
结果表明，BEV与BEV检测器与仅BEV检测器表现相当，远不及透视与BEV检测器。
因此，仅在透视图上构建第一阶段头并施加辅助监督对BEV模型是有帮助的。

%\vspace{-0.3cm}
\subsubsection{透视监督的泛化性}
本文提出的透视监督预期适用于不同架构和大小的骨干网络。
本文在一系列常用于3D目标检测任务的骨干网络上构建 BEVFormer v2：ResNet \cite{ResNet}、DLA \cite{DLA}、VoVNet \cite{Vovnet} 和 InternImage \cite{InternImage}。结果报告在表 \ref{table:large_backbone} 中。
与纯BEV检测器相比，BEVFormer v2（BEV与透视）在所有骨干网络上的NDS提升约3\%，mAP提升约2\%，表明其适用于不同架构和模型大小。
本文认为额外的透视监督可以成为训练BEV模型的通用方案，特别是在没有3D预训练的情况下适配大规模图像骨干网络时。

\setlength{\tabcolsep}{3pt}
\setlength{\doublerulesep}{2\arrayrulewidth}
\renewcommand{\arraystretch}{1.1}
\begin{table*}[ht]
    \caption{BEVFormer v2 中透视图头（Perspective Head）与 BEV 头（BEV Head）不同选择的对比。模型在 nuScenes 验证集上评测。所有模型均未使用时序（Temporal）信息训练。}
    \label{table:detector_choice_for_two_view}
    \centering
    \vspace{-5pt}
    \begin{tabular}{cc|c|c|cc|ccccc}
        \toprule
        Perspective View & BEV View & Backbone & Epoch & NDS & mAP & mATE & mASE & mAOE & mAVE & mAAE   \\ 
        \midrule 
        DD3D & Deformable DETR & ResNet-50 & 48 & 0.428 & 0.349 & 0.750 & 0.276 & 0.424 & 0.817 & 0.193 \\ 
        DD3D & Group DETR & ResNet-50 & 48 & 0.445 & 0.353 & 0.725 & 0.276 & 0.366 & 0.767 & 0.180 \\ 
         DETR3D & Deformable DETR & ResNet-50 & 48 & 0.409 & 0.335 & 0.765 & 0.276 & 0.469 & 0.877 & 0.198 \\ 
        DETR3D & Group DETR & ResNet-50 & 48 & 0.423 & 0.351 & 0.743 & 0.279 & 0.466 & 0.844 & 0.201 \\ 
        \midrule
        \bottomrule
    \end{tabular}
    
\end{table*}
\let\ck\checkmark

\setlength{\tabcolsep}{3pt}
\setlength{\doublerulesep}{2\arrayrulewidth}
\renewcommand{\arraystretch}{1.0}
\begin{table*}[t]
    \caption{BEVFormer v2 各组件（Bells and Whistles）的消融实验（Ablation Study），在 nuScenes 验证集上评测。所有模型均使用 ResNet-50 骨干网络（Backbone）及时序（Temporal）信息训练。`Pers'、`IDA'、`Long' 和 `Bi' 分别表示透视监督（Perspective Supervision）、图像级数据增强（Image-level Data Augmentation）、长时序间隔（Long Temporal Interval）和双向时序编码器（Bi-directional Temporal Encoder）。}
    \label{table:trick}
    \centering
    \vspace{-5pt}
    \begin{tabular}{l|c|cccc|cc|cc|cccccc}
        \toprule
        Method & Epoch & Pers & IDA & Long & Bi & NDS & mAP & mATE & mASE & mAOE & mAVE & mAAE   \\ 
        \midrule
        Baseline & 24 & \ck & & & & 0.478 & 0.368 & 0.709 & 0.282 & 0.452 & 0.427 & 0.191 \\ 
        Image-level Data Agumentation & 24 & \ck & \ck & & & 0.489 & 0.386 & 0.690 & 0.273 & 0.482 & 0.395 & 0.199 \\
        % BDA  & & & & & & & \\
        Longer Temporal Interval & 24 & \ck & \ck & \ck & & 0.498 & 0.388 & 0.679 & 0.276 & 0.417 & 0.403 & 0.189 \\
        % Randomized Temporal   & \ck & \ck & \ck & \ck & & & \\ 
        Bi-directional Temporal Encoder & 24 & \ck & \ck & \ck & \ck & 0.529 & 0.423 & 0.618 & 0.273 & 0.413 & 0.333 & 0.181 \\
        All but Perspective & 24 & & \ck & \ck & \ck & 0.507 & 0.397 & 0.636 & 0.281 & 0.455 & 0.356 & 0.190 \\
        \bottomrule
    \end{tabular}
    
\end{table*}

\subsubsection{训练轮数的选择}
本文对仅BEV模型和 BEVFormer v2（BEV与透视）进行了不同轮数的训练，观察两个模型达到收敛所需的时间。
表 \ref{table:long_schedule} 显示，本文的BEV与透视模型比仅BEV模型收敛更快，证实了辅助透视损失促进了优化。
仅BEV模型在更长时间的训练下仅获得边际提升。
但在72个epoch时两个模型之间的差距仍然存在，且即使更长时间的训练也可能无法消除，这表明仅靠BEV监督无法很好地适配图像骨干网络。
根据表 \ref{table:long_schedule}，本文模型训练48个epoch就已足够，其他实验除非另有说明均保持此设置。

\subsubsection{检测头的选择}
BEVFormer v2 可使用多种类型的透视和BEV检测头。
本文探索了几种代表性方法以选择最佳组合：对于透视头，候选方案为 DD3D \cite{DD3D} 和 DETR3D \cite{DETR3D}；对于BEV头，候选方案为可变形 DETR \cite{deformable-detr} 和 Group DETR \cite{groupdetr}。
DD3D 是一个单阶段无锚框（Anchor-free）透视头，在图像特征上进行密集的逐像素预测。
相反，DETR3D 使用3D到2D查询来采样图像特征并提出稀疏的预测集。
然而，根据本文的定义，它属于透视监督，因为它利用图像特征进行最终预测而不生成BEV特征，即损失直接作用于图像特征。
如表 \ref{table:detector_choice_for_two_view} 所示，DD3D 作为透视头优于 DETR3D，支持了本文在章节 \ref{sec:perspective_supervision} 中的分析。
DD3D 提供的密集直接监督对BEV模型有帮助，而 DETR3D 的稀疏监督无法克服BEV头的缺陷。
Group DETR 头是可变形 DETR 头的扩展，使用分组物体查询和组内自注意力。
Group DETR 在BEV头上取得了更好的性能，但计算量更大。
因此，本文在表 \ref{table:sota_nuscenes} 中采用了 DD3D 头和 Group DETR 头，在其他消融实验中保持与 BEVFormer \cite{bevformer} 相同的可变形 DETR 头。

%\vspace{-10pt}
\subsubsection{各组件的消融实验}
\label{sec:bells_and_whistles}
在表 \ref{table:trick} 中，本文对 BEVFormer v2 中采用的各种组件进行了消融实验，以确认其对最终结果的贡献，包括：
(1) 图像级数据增强（Image-level Data Augmentation, IDA）——图像随机水平翻转。
(2) 更长时序间隔——BEVFormer \cite{bevformer} 使用0.5秒间隔的连续帧，而本文的 BEVFormer v2 以2秒间隔采样历史BEV特征。
(3) 双向时序编码器（Bi-directional Temporal Encoder）——对于离线3D检测，本文 BEVFormer v2 中的时序编码器可以利用未来的BEV特征。
通过更长的时序间隔，模型可以从更多不同时间戳的自车位置收集信息，有助于估计物体朝向并显著降低mAOE。
在离线3D检测设置中，双向时序编码器能够提供来自未来帧的额外信息，大幅提升模型性能。
本文还在应用所有组件的情况下对透视监督进行了消融实验。
如表 \ref{table:trick} 所示，透视监督将NDS提升了2.2\%，mAP提升了2.6\%，是主要改进来源。
